\begin{tikzpicture}[wf]
  \draw[wf frame] (0,0) rectangle (13,20);

  \draw[wf bar] (0,19.1) rectangle (13,20);
  \draw[black!70, line width=0.7pt] (0.35,19.35) -- (0.35,19.65) -- (0.55,19.85)
      -- (0.75,19.65) -- (0.75,19.35) -- cycle;
  \node[wf btn] at (2.65,19.55) {Számításelmélet I $\vee$};
  \node[wf btn] at (6.45,19.55) {Számításelmélet II $\vee$};

  \node[anchor=west, font=\sffamily\scriptsize\itshape\bfseries, black!60] at (1,18.5)
      {Számításelmélet I > Reguláris nyelvek > Determinisztikus véges automata};
  \node[wf mark] at (12.5,18.5) {1};
  \node[font=\sffamily\Large\bfseries] at (6.5,17.6) {Determinisztikus véges automata};

  \draw[wf box] (1,14.6) rectangle (12,16.9);
  \node[anchor=north west, font=\sffamily\footnotesize\bfseries] at (1.15,16.85) {Definíció};
  \foreach \y/\l in {15.95/11.4, 15.5/10.8, 15.05/7.5}
    \draw[wf line] (1.4,\y) -- (\l,\y);
  \draw[wf box] (1,12.6) rectangle (12,14.3);
  \node[anchor=north west, font=\sffamily\footnotesize\bfseries] at (1.15,14.25) {Tétel};
  \foreach \y/\l in {13.4/11.2, 12.95/8.3}
    \draw[wf line] (1.4,\y) -- (\l,\y);
  \node at (6.5,12.2) {$\delta(q, a) = q'$};
  \node[anchor=west, font=\sffamily\footnotesize\itshape, black!60] at (1,11.75) {\ldots};
  \node[wf mark] at (12.5,16.9) {2};

  \node[anchor=west, font=\sffamily\large\bfseries] at (1,11.0) {Próbáld ki!};
  \node[wf mark] at (12.5,11.0) {3};
  \node[anchor=west] at (1,10.25) {Bemeneti szó:};
  \node[wf input, minimum width=3.6cm] at (3.4,10.25) {\textcolor{black!40}{$\varepsilon$ (üres szó)}};
  \foreach \s [count=\i] in {{$|\!\leftarrow$}, {$\leftarrow$}, {$\triangleright$}, {$\rightarrow$}, {$\rightarrow\!|$}}
    \node[wf btn, minimum width=0.7cm] at (7.4+0.8*\i,10.25) {\s};
  \draw[black!50, line width=1.5pt] (8.0,9.85) -- (12,9.85);
  \fill[black!70] (9.2,9.85) circle (0.09);

  \draw[wf box] (1,6.6) rectangle (12,9.6);
  \node[wf state, fill=black!20, line width=1.2pt] (q0) at (3,8.1) {$q_0$};
  \node[wf state] (q1) at (6.5,8.1) {$q_1$};
  \node[wf state, double, double distance=1.2pt] (q2) at (10,8.1) {$q_2$};
  \draw[->, >=stealth] (1.8,8.1) -- (q0);
  \draw[->, >=stealth] (q0) -- node[above] {$a$} (q1);
  \draw[->, >=stealth] (q1) -- node[above] {$b$} (q2);
  \draw[->, >=stealth] (q0) to[loop below] node[right] {$b$} (q0);
  \draw[->, >=stealth] (q2) to[out=-140, in=-40] node[below] {$a, b$} (q1);
  \node[anchor=west] at (1,6.1) {Konfigurációk: $(q_0, ab) \vdash (q_1, b) \vdash \ldots$};

  \node[anchor=west, font=\sffamily\large\bfseries] at (1,5.0) {Teszteld magad!};
  \node[wf mark] at (12.5,5.0) {4};
  \node[anchor=west, font=\sffamily\footnotesize\bfseries] at (1,4.45) {Könnyű};
  \draw[wf box] (1,1.55) rectangle (12,4.1);
  \node[anchor=west] at (1.3,3.7) {\textbf{1.} Elfogadja-e az automata az \emph{abba} szót?};
  \draw[draw=black!50, fill=black!15, rounded corners=2pt] (1.3,2.75) rectangle (11.7,3.25);
  \node[anchor=west] at (1.45,3.0) {\textbf{A)} Igen};
  \draw[draw=black!50, rounded corners=2pt] (1.3,2.1) rectangle (11.7,2.6);
  \node[anchor=west] at (1.45,2.35) {\textbf{B)} Nem};
  \node[anchor=west, font=\sffamily\scriptsize] at (1.3,1.8) {\textbf{Helyes!} Magyarázat \ldots};

  \draw[wf bar] (0,0) rectangle (13,1.2);
  \node[anchor=west, align=left] at (0.5,0.6)
      {Készítette: Szabó Anna\\ ELTE Informatikai Kar $\cdot$ 2026};
\end{tikzpicture}
